\chapter{Inverse and Implicit Function Theorems}
\label{sec:inverse-function-theorem}
\noindent\emph{Editorial supplement. This appendix is not part of Warner's original text.}
\medskip

In this chapter we discuss three logically equivalent theorems: the inverse
function theorem, the implicit function theorem, and the constant rank
theorem, which describe the local behavior of $C^\infty$ maps from $\RR^n$ to
$\RR^m$. We assume the inverse function theorem and then derive the other two
theorems from it in the simplest cases.


In subsequent discussions, these theorems are applied on manifolds to study
the local behavior of $C^\infty$ maps, where the maps in question usually have
maximal rank at a point or constant rank in a neighborhood.


\section{The Inverse Function Theorem}

A $C^\infty$ map $f: U\to \RR^n$ defined on an open set $U\subset\RR^n$ is
said to be \emph{locally invertible}, or to be a \emph{local diffeomorphism}
near a point $p\in U$, if $f$ has a $C^\infty$ inverse map in some
neighborhood of $p$. The inverse function theorem gives a criterion for a map
to be locally invertible. We call the matrix $Jf = [\partial f^i/\partial
	x^j]$ the \emph{Jacobian matrix} of $f$, and its determinant $\det
	[\partial f^i/\partial x^j]$ is called the \emph{Jacobian determinant of
	$f$}.

\begin{theorem}[The Inverse Function Theorem]
  \label{thm:inverse-function-theorem-Rn}

	Let $f:U\to\RR^n$ be a $C^\infty$ map defined on an open subset
	$U\subset\RR^n$. At any point $p\in U$, the map $f$ is
	\underline{locally invertible} near $p$ \underline{if and only if} the
	determinant of its Jacobian matrix $\det [\partial f^i/\partial x^j
			(p)]$ is nonzero.

\end{theorem}

% A proof of the theorem can be found in \cite{rudinma}, p. 221, Theorem 9.24.
Note that although the inverse function theorem requires only the Jacobian
determinant at the point $p$ to be nonzero, continuity in fact guarantees
that the Jacobian determinant is nonzero throughout a neighborhood of $p$.


Since the linear map defined by the Jacobian matrix $Jf(p)$ is the best
approximation to $f$ at $p$, it is quite plausible that $f$ being invertible
in a neighborhood of $p$ is equivalent to $Jf(p)$ being invertible (or,
equivalently, to $\det Jf(p)\neq 0$).


\begin{remark}
	Theorem \ref{thm:inverse-function-theorem-Rn} gives a necessary and
	sufficient condition. However, if one only requires the map $f$ to be a
	local homeomorphism (without regard to smoothness), the theorem gives
	only a sufficient condition.
\end{remark}




\section{The Implicit Function Theorem}

As an application of the inverse function theorem, we state the following
implicit function theorem. Let $F:U\subset\RR^m\to\RR^n$ be a $C^\infty$ map,
where $m>n$. If $p\in U$ is a solution of the equation $F(x^1,\cdots,
	x^m)=\bm{0}$, we would like to ask: in some neighborhood of this
solution, do there exist other solutions? If so, how can these solutions be
represented? The implicit function theorem answers these questions in terms of
the rank of the Jacobian matrix of $F$ at $p$.



\begin{theorem}[Implicit Function Theorem]
\label{thm:appendix-implicit}
Let $F$ be a smooth map from an open neighborhood of $(a,b)$ in
$\RR^r\times\RR^n$ to $\RR^n$, with $F(a,b)=0$. If the partial derivative
$D_yF(a,b):\RR^n\to\RR^n$ is invertible, then there are neighborhoods
$A$ of $a$ and $B$ of $b$ and a unique smooth map $g:A\to B$ such that
\[
 F(x,y)=0\quad\Longleftrightarrow\quad y=g(x)
 \qquad((x,y)\in A\times B).
\]
Moreover, $g(a)=b$ and
\[
 Dg(x)=-[D_yF(x,g(x))]^{-1}D_xF(x,g(x)).
\]
\end{theorem}
\begin{proof}
Define $\Phi(x,y)=(x,F(x,y))$. Its derivative at $(a,b)$ is the block matrix
\[
 D\Phi(a,b)=\begin{pmatrix}I&0\\D_xF(a,b)&D_yF(a,b)\end{pmatrix},
\]
which is invertible. By the inverse function theorem, $\Phi$ is a
diffeomorphism between neighborhoods of $(a,b)$ and $(a,0)$. Its inverse
has the form $(u,v)\mapsto(u,h(u,v))$ because $\Phi$ preserves the first
coordinate. Set $g(u)=h(u,0)$ and shrink to product neighborhoods so that
the graph remains in the domain of this inverse. Injectivity of $\Phi$
gives uniqueness and the displayed equivalence. Differentiating
$F(x,g(x))=0$ gives the derivative formula, after shrinking the neighborhoods
so that $D_yF$ remains invertible.
\end{proof}

\begin{corollary}[Regular Level Sets]
If $F:U\subset\RR^m\to\RR^n$ has rank $n$ at $p$ and $F(p)=c$, then
near $p$ the level set $F^{-1}(c)$ is the graph of a smooth function in
$m-n$ variables. In particular, it is a submanifold of dimension $m-n$.
\end{corollary}
\begin{proof}
A nonzero $n\times n$ minor permits a permutation of the variables for
which $D_yF(p)$ is invertible. Apply the implicit function theorem to $F-c$.
\end{proof}

\section{The Constant Rank Theorem}
\begin{theorem}[Constant Rank Theorem]
\label{thm:appendix-constant-rank}
Let $f:U\subset\RR^m\to\RR^n$ be smooth and have constant rank $k$ in
a neighborhood of $p$. There are smooth coordinates centered at $p$ and
$f(p)$ in which $f$ has the form
\[
 (u_1,\ldots,u_m)\longmapsto(u_1,\ldots,u_k,0,\ldots,0).
\]
\end{theorem}
\begin{proof}
For $k=0$, all first derivatives vanish near $p$, so $f$ is constant on a
sufficiently small ball and the assertion follows by translating the target.
For $k>0$, translate $p$ and $f(p)$ to zero and permute the coordinates so
that $\partial(f_1,\ldots,f_k)/\partial(x_1,\ldots,x_k)$ is invertible at
zero. The inverse function theorem applies to
\[
 \Phi(x)=(f_1(x),\ldots,f_k(x),x_{k+1},\ldots,x_m).
\]
In the coordinates $u=\Phi(x)$, the map is
$f\circ\Phi^{-1}(u)=(u_1,\ldots,u_k,h_{k+1}(u),\ldots,h_n(u))$.
Its derivative has the block form
\[
 \begin{pmatrix}I_k&0\\ A&B\end{pmatrix}.
\]
Since its rank is exactly $k$ throughout a neighborhood, $B=0$ there.
On a smaller coordinate box, each $h_j$ therefore depends only on
$u_1,\ldots,u_k$. The target coordinate change
\[
 \Psi(y)=\bigl(y_1,\ldots,y_k,
 y_{k+1}-h_{k+1}(y_1,\ldots,y_k),\ldots,
 y_n-h_n(y_1,\ldots,y_k)\bigr)
\]
has a smooth inverse obtained by adding the same functions. Thus
$\Psi\circ f\circ\Phi^{-1}$ has the asserted form.
\end{proof}

When $k=m$, this is the local form of an immersion; when $k=n$, it is
the local form of a submersion. Unlike the regular-level-set case, a rank
condition at just one point is not sufficient for the general constant-rank
conclusion. For example, $x\mapsto x^2$ has rank zero at zero but is not
locally constant.

\section{Supplementary Linear Algebra}

\subsection{Dual Spaces}\label{sec:dualspace}

If $V$ and $W$ are real linear spaces, we denote by $\Hom(V,W)$ the vector
space of all linear maps $f: V\to W$. The \emph{dual space} $V^\vee$ of $V$ is
defined to be the linear space of all linear real-valued functions on $V$:
\begin{displaymath}
	V^\vee = \Hom(V,\mathbb{R}).
\end{displaymath}
The elements of $V^\vee$ are called \emph{covectors} or $1$-covectors on $V$.

In the discussion below, we always assume that $V$ is a \emph{finite
	dimensional} vector space. Let $e_1,\cdots, e_n$ be a basis of $V$. Then
every vector $v\in V$ can be uniquely expressed as a linear combination $v=
	\sum v^i e_i$, where $v^i\in\mathbb{R}$. Let $\alpha^i: V\to \mathbb{R}$
be the linear function that takes the value of the $i$-th coordinate component
of $v$, $\alpha^i(v) = v^i$. Then
\begin{displaymath}
	\alpha^i(e_j) = \delta^i_j =
	\begin{cases}
		1 & \text{if $i=j$},     \\
		0 & \text{if $i\neq j$}.
	\end{cases}
\end{displaymath}

\begin{proposition}
	The functions $\alpha^1,\cdots, \alpha^n$ form a basis of $V^\vee$.
\end{proposition}

\begin{proof}
For $\lambda\in V^\vee$ and $v=\sum_i v^ie_i$, linearity gives
$\lambda(v)=\sum_i\lambda(e_i)v^i$. Hence
$\lambda=\sum_i\lambda(e_i)\alpha^i$, so the $\alpha^i$ span $V^\vee$.
Evaluating a relation $\sum_i c_i\alpha^i=0$ on each $e_j$ gives $c_j=0$,
which proves linear independence.
\end{proof}


\begin{corollary}
	The dual space $V^\vee$ has the same dimension as the finite dimensional
	vector space $V$.
\end{corollary}

\subsection{Dual Maps and Annihilators}
For a linear map $T:V\to W$, its dual map is
$T^*:W^\vee\to V^\vee$, $T^*(\lambda)=\lambda\circ T$.
For a subspace $E\subset V$, its annihilator is
$E^\circ=\{\lambda\in V^\vee:\lambda|_E=0\}$.
Extending a basis of $E$ to a basis of $V$ shows that
$\dim E^\circ=\dim V-\dim E$.

\begin{proposition}
For finite-dimensional $V$ and $W$,
\[
 \ker T^*=(\operatorname{im}T)^\circ,\qquad
 \operatorname{im}T^*=(\ker T)^\circ.
\]
In particular, $\operatorname{rank}T^*=\operatorname{rank}T$.
\end{proposition}
\begin{proof}
The first equality follows directly from $T^*\lambda=\lambda\circ T$.
Every element of $\operatorname{im}T^*$ vanishes on $\ker T$, giving one
inclusion in the second equality. Conversely, if $\mu$ vanishes on
$\ker T$, the rule $Tv\mapsto\mu(v)$ defines a linear functional on
$\operatorname{im}T$. Extend it to $W$ using a basis; the extension $\lambda$
satisfies $T^*\lambda=\mu$. The rank identity now follows from the
dimension formula for annihilators and rank--nullity.
\end{proof}
